\section{Contextualização e objetivo}
\label{sec:contexto}
Este documento compara os tipos societários previstos no ordenamento jurídico nacional, fundamenta a estrutura jurídica proposta para a startup e identifica aspectos legais relevantes para sua formalização.

O \projeto{} é uma plataforma digital voltada à prevenção e ao acesso a cuidados para animais de companhia. Seu modelo reúne lembretes e registros do pet, orientação e encaminhamento, clínicas parceiras, tabela social e créditos solidários. A proposta contempla receitas de assinaturas, planos para parceiros, intermediação de atendimentos e convênios ou patrocínios. O lançamento previsto é local, em Belo Horizonte, com expansão futura.

\subsection{Premissas da análise}
\begin{itemize}
\item \textbf{Equipe:} quatro integrantes, identificados na capa e na folha de rosto. Para esta análise, propõe-se que todos sejam sócios fundadores. A condição de integrante acadêmico não comprova participação societária efetiva.
\item \textbf{Fase:} projeto acadêmico em desenvolvimento, com projeções e hipóteses de operação. Não se presume a existência de CNPJ, registro de marca, licenças ou contratos já celebrados.
\item \textbf{Atuação:} plataforma de tecnologia e conexão com prestadores; não uma seguradora, plano de saúde ou clínica própria. A execução efetiva do serviço, e não apenas sua denominação, determinará as exigências regulatórias.
\item \textbf{Decisões propostas:} LTDA com quatro sócios, avaliação inicial do Simples Nacional e participação igualitária ilustrativa. Essas escolhas dependem de confirmação do grupo e de validação contábil e jurídica antes da abertura.
\end{itemize}

\subsection{Forma jurídica, porte e tributação}
Essas categorias não são equivalentes. \textbf{LTDA} é uma forma societária; \textbf{ME} e \textbf{EPP} são enquadramentos de porte; \textbf{Simples Nacional} é um regime tributário. O \textbf{MEI} é um enquadramento simplificado do empresário individual, não uma sociedade com sócios. A \textbf{SLU} é uma limitada com um único sócio, e não um tipo societário independente da LTDA (Código Civil, 2002; LC nº 123/2006).

A designação \emph{startup} também não substitui a forma jurídica. O enquadramento no Marco Legal das Startups exige os requisitos da LC nº 182/2021, incluindo receita bruta de até R\$~16 milhões no ano anterior, até dez anos de inscrição no CNPJ e atendimento às condições de inovação previstas na lei. Não decorre automaticamente do uso de um aplicativo (LC nº 182/2021).

\textbf{Escopo.} As conclusões são acadêmicas, não um parecer profissional nem comprovação de formalização. As referências indicam as normas utilizadas; sua versão consolidada e os procedimentos administrativos devem ser conferidos na data da abertura.